\scanpage{64}
Que \unsure{deviendraient|donneraient} les conjectures précédentes, quand on les applique \unsure{à|:} une situation où $K$ est remplacé par un modèle $S$ de $K$ (disons, élémentaire anabélien), quand $U_K$ provient d'une courbe relative\unsure{.} $U_S$ sur $S$ — de sorte qu'on a un hom de groupes
\diagram{\begin{tikzcd}[ampersand replacement=\&,row sep=1.2em,column sep=6em]
{} \& E_{U_S}\arrow[dd] \\
(4) \& {} \\
{} \& E_S
\end{tikzcd}}
de noyau $\pi_{\bar U}$, dont $E_{U_K}-E_K$ est déduit par chgt de base i.e. par produit fibré
\diagram{\begin{tikzcd}[ampersand replacement=\&,row sep=1.2em,column sep=4em]
{} \& \pi\arrow[r,no head]\arrow[d] \& \pi\arrow[d] \\
{} \& E_{U_S}\arrow[dd] \& E_{U_K}\arrow[l]\arrow[dd] \\
\unsure{(5)|(4')} \& {} \& {} \\
{} \& E_S \& E_K\arrow[l,"\mathrm{surj}"']\arrow[uul,dashed,"\varphi"']
\end{tikzcd}}
Ainsi, les sections de $E_{U_{\deleted{\unsure{S}}K}}$ sur $E_{\deleted{\unsure{S}}K}$ correspondent aux \deleted{\uline{\unsure{fausses|bonnes}} sections, de} \inserted{relèvements continus} $E_K\longrightarrow E_{U_S}$ de l'hom surjectif $E_K\longrightarrow E_S$, et parmi ces relèvements, ceux qui sont \deleted{\unsure{nuls}} \inserted{triviaux} sur le noyau de $E_K\longrightarrow E_S$ correspondent exactement aux sections
\margin{63}

